\section{Frequency Analysis}

Frequency analysis mengeksploitasi fakta bahwa huruf dalam bahasa alami memiliki frekuensi kemunculan yang dapat diprediksi \cite{trappe2006,schneier1996}.

\begin{table}[!htbp]
\centering
\caption{Frekuensi huruf dalam bahasa Inggris (\%)}
\label{tab:frequency}
\begin{tabular}{@{}*{7}{c}@{}}
\toprule
E 12.7 & T 9.1 & A 8.2 & O 7.5 & I 7.0 & N 6.7 & S 6.3 \\
H 6.1 & R 6.0 & D 4.3 & L 4.0 & C 2.8 & U 2.8 & M 2.4 \\
W 2.4 & F 2.2 & G 2.0 & Y 2.0 & P 1.9 & B 1.5 & V 1.0 \\
K 0.8 & J 0.2 & X 0.2 & Q 0.1 & Z 0.1 & & \\
\bottomrule
\end{tabular}
\end{table}

Dalam bahasa Inggris, E, T, A, O, I, N paling sering muncul; J, Q, X, Z paling jarang.

Untuk cipher substitusi monoalfabetik: hitung frekuensi setiap huruf dalam ciphertext, bandingkan dengan frekuensi standar bahasa, dan cocokkan pemetaan \cite{trappe2006,schneier1996}. Ciphertext yang lebih panjang memberikan estimasi frekuensi lebih akurat.
